%% =========================================================================
%%  cbf_logdet_triangulation_multi_feature_revised.tex
%%  Control Barrier Function for Multi-Feature Triangulation Observability
%%  using an Average Log-Determinant Safe Set
%% =========================================================================

\documentclass[11pt,a4paper]{article}

%% ---------- Packages ----------
\usepackage[margin=2.5cm]{geometry}
\usepackage{amsmath,amssymb,amsthm,mathtools}
\usepackage{bm}
\usepackage{hyperref}
\usepackage{cleveref}
\usepackage{booktabs}
\usepackage{xcolor}
\usepackage{mdframed}
\usepackage{thmtools}
\usepackage{enumitem}
\usepackage{float}
\usepackage{algorithm}
\usepackage{algpseudocode}
\usepackage{cancel}

%% ---------- Theorem environments ----------
\declaretheoremstyle[
  spaceabove=6pt, spacebelow=6pt,
  headfont=\normalfont\bfseries,
  notefont=\normalfont\bfseries, notebraces={(}{)},
  bodyfont=\normalfont,
  postheadspace=5pt
]{mystyle}

\declaretheorem[style=mystyle,name=Theorem,numberwithin=section]{theorem}
\declaretheorem[style=mystyle,name=Lemma,numberwithin=section]{lemma}
\declaretheorem[style=mystyle,name=Proposition,numberwithin=section]{proposition}
\declaretheorem[style=mystyle,name=Definition,numberwithin=section]{definition}
\declaretheorem[style=mystyle,name=Remark,numberwithin=section,shaded={bgcolor=gray!10}]{remark}

%% ---------- Coloured result boxes ----------
\newmdenv[
  backgroundcolor=blue!6,
  linecolor=blue!40,
  linewidth=1pt,
  roundcorner=4pt,
  skipabove=8pt,
  skipbelow=8pt
]{resultbox}

\newmdenv[
  backgroundcolor=orange!8,
  linecolor=orange!60,
  linewidth=1pt,
  roundcorner=4pt,
  skipabove=8pt,
  skipbelow=8pt
]{keyresult}

%% ---------- Macros ----------
\newcommand{\R}{\mathbb{R}}
\newcommand{\SO}{\mathrm{SO}(3)}
\newcommand{\Stwo}{\mathcal{S}^{2}}
\newcommand{\bzero}{\bm{0}}
\newcommand{\bv}{\bm{v}}
\newcommand{\bo}{\bm{\omega}}
\newcommand{\bu}{\bm{u}}
\newcommand{\bb}{\bm{b}}
\newcommand{\bp}{\bm{p}}
\newcommand{\bR}{\bm{R}}
\newcommand{\bI}{\bm{I}}
\newcommand{\bM}{\bm{M}}
\newcommand{\bPsi}{\bm{\Psi}}
\newcommand{\sx}[1]{[#1]_{\times}}
\newcommand{\logdet}{\log\det}
\newcommand{\Tr}{\operatorname{tr}}
\newcommand{\bij}{\bm{b}_{i,j}}
\newcommand{\bhij}{\hat{\bm{b}}_{i,j}}
\newcommand{\pij}{\bm{\pi}_{i,j}}
\newcommand{\Mii}{\bm{M}_{i}}
\newcommand{\ellii}{\ell_i}

%% ---------- Hyperref ----------
\hypersetup{
  colorlinks=true,
  linkcolor=blue!60!black,
  citecolor=green!50!black,
  urlcolor=blue!80!black
}

%% =========================================================================
\begin{document}

\title{\textbf{Control Barrier Function for Multi-Feature Triangulation Observability}\\[4pt]
       \large An Average Log-Determinant Safe-Set Formulation}
\author{}
\date{\today}
\maketitle
\tableofcontents
\bigskip

% =========================================================================
\section{Introduction and Problem Motivation}
% =========================================================================

In vision-aided navigation and visual-inertial odometry (VIO), the quality of
triangulated 3D feature positions is a fundamental prerequisite for
state-estimation accuracy.  Triangulation becomes \emph{ill-conditioned}---or
even \emph{impossible}---when a camera undergoes pure rotation without
translation, or when the observed bearing vectors provide insufficient
parallax.  These scenarios cause the triangulation information matrix to
become rank deficient or poorly conditioned.

This document considers a finite trajectory window containing $n$ camera poses
and $s$ visible or selected point features.  The index
$i\in\{1,\ldots,s\}$ denotes the feature index, while
$j\in\{1,\ldots,n\}$ denotes the camera pose inside the trajectory window.
Thus, the bearing associated with feature $i$ observed from pose $j$ is denoted
by $\bm{b}_{i,j}$.

\begin{definition}[Per-feature triangulation observability]
For feature $i$, the set of bearing measurements
$\{\bm{b}_{i,j}\}_{j=1}^{n}$ admits a unique least-squares triangulation if
and only if the corresponding information matrix
\begin{equation}
  \bm{M}_{i}
  = \sum_{j=1}^{n}\left(\eta_{i,j}^{2}\bm{I}
  - \bm{b}_{i,j}\bm{b}_{i,j}^{\top}\right)
  \label{eq:intro_Mi}
\end{equation}
is positive definite, where $\eta_{i,j}=\|\bm{b}_{i,j}\|_2$.  If the bearings
are unit-normalized, then $\eta_{i,j}=1$ and
$\bm{M}_{i}=\sum_{j=1}^{n}(\bm{I}-\bm{b}_{i,j}\bm{b}_{i,j}^{\top})$.
\end{definition}

The goal is to derive a \textbf{Control Barrier Function (CBF)} that modifies
the nominal camera translational velocity command so that the average
triangulation quality over all selected features remains above a prescribed
threshold.  Each feature has its own log-determinant observability metric,
\begin{equation}
  \ell_i = \log\det\bm{M}_i,
\end{equation}
and the CBF reads the average value
\begin{equation}
  \bar{\ell}
  = \frac{1}{s}\sum_{i=1}^{s}\ell_i
  = \frac{1}{s}\sum_{i=1}^{s}\log\det\bm{M}_i.
  \label{eq:average_logdet_intro}
\end{equation}
The safety objective is therefore $\bar{\ell}\ge h_{\min}$, where $h_{\min}$
is a user-selected observability threshold.

% =========================================================================
\section{Notation and Conventions}
% =========================================================================

\begin{table}[h]
\centering
\renewcommand{\arraystretch}{1.25}
\begin{tabular}{clc}
\toprule
\textbf{Symbol} & \textbf{Description} & \textbf{Space} \\
\midrule
$i\in\{1,\ldots,s\}$ & Feature index & $\mathbb{N}$ \\
$j\in\{1,\ldots,n\}$ & Camera-pose index in the trajectory window & $\mathbb{N}$ \\
${}^{A}\bm{p}_{f_i}$ & Position of feature $i$ in anchor frame $A$ & $\R^3$ \\
${}^{A}\bm{p}_{C_j}$ & Position of camera pose $C_j$ in anchor frame $A$ & $\R^3$ \\
$\bm{R}_{C_j}^{A}$ & Rotation matrix from camera frame $C_j$ to anchor frame $A$ & $\SO$ \\
$\bm{R}_{B}^{C}$ & Rotation matrix from body frame $B$ to camera frame $C$ & $\SO$ \\
$\bm{v}_{B,j}$ & Body-frame velocity at pose $j$ & $\R^3$ \\
$\bm{v}_{B}$ & Current velocity command $\bm{v}_{B,n}$ (control input) & $\R^3$ \\
$f(\bm{x})$ & Drift term (past-pose contribution to $\dot{h}$) & $\R$ \\
$\bm{g}(\bm{x})$ & Control-input gradient vector in body frame & $\R^3$ \\
$\bm{b}_{i,j}$ & Bearing of feature $i$ from pose $j$, expressed in anchor frame & $\R^3$ \\
$\hat{\bm{b}}_{i,j}$ & Unit bearing $\bm{b}_{i,j}/\eta_{i,j}$ & $\mathcal{S}^2$ \\
$\rho_{i,j}$ & Depth/range of feature $i$ from pose $j$ & $\R_{>0}$ \\
$\bm{\pi}_{i,j}$ & Projection/information contribution of bearing $\bm{b}_{i,j}$ & $\R^{3\times3}$ \\
$\bm{M}_{i}$ & Triangulation information matrix for feature $i$ & $\R^{3\times3}$ \\
$\ell_i$ & Feature-wise log-det metric $\log\det\bm{M}_i$ & $\R$ \\
$\bar{\ell}$ & Average log-det metric over all selected features & $\R$ \\
\bottomrule
\end{tabular}
\caption{Summary of the multi-feature notation.}
\label{tab:notation}
\end{table}

The notation ${}^{A}\bm{p}_{C_j}$ is used throughout the document instead of
${}^{A}\bm{p}_{C}$ because $j$ explicitly denotes the camera pose inside the
trajectory window.  Similarly, any pose-dependent rotation from the camera
frame to the anchor frame is written as $\bm{R}_{C_j}^{A}$.

% =========================================================================
\section{System Model and Feature Geometry}
% =========================================================================

For feature $i$ observed from camera pose $j$, the position of the feature in
camera frame $C_j$ is
\begin{equation}
  {}^{C_j}\bm{p}_{f_i}
  = \bm{R}_{A}^{C_j}\left({}^{A}\bm{p}_{f_i}
  - {}^{A}\bm{p}_{C_j}\right)
  \in\R^3 .
  \label{eq:feature_camera_frame}
\end{equation}
Equivalently, after rotating this vector back to the anchor frame,
\begin{equation}
  {}^{C_j\to A}\bm{p}_{f_i}
  = \bm{R}_{C_j}^{A}\,{}^{C_j}\bm{p}_{f_i}
  = {}^{A}\bm{p}_{f_i}-{}^{A}\bm{p}_{C_j}
  = \rho_{i,j}\,\bm{b}_{i,j}.
  \label{eq:feature_anchor_factored}
\end{equation}
Hereafter, $\bm{b}_{i,j}$ denotes the bearing of feature $i$ observed from pose
$j$ and expressed in the anchor frame.  If the raw bearing is not unit length,
then
\begin{equation}
  \bm{b}_{i,j}=\eta_{i,j}\hat{\bm{b}}_{i,j},
  \qquad
  \eta_{i,j}=\|\bm{b}_{i,j}\|_2,
  \qquad
  \|\hat{\bm{b}}_{i,j}\|_2=1 .
\end{equation}

The pose-dependent translational kinematics are written as
\begin{equation}
  {}^{A}\dot{\bm{p}}_{C_j}
  = \bm{R}_{C_j}^{A}\bm{R}_{B}^{C}\bm{v}_{B,j}.
  \label{eq:pose_velocity}
\end{equation}
Here $\bm{v}_{B,j}$ denotes the body-frame velocity at pose $j$.

\begin{remark}[Past vs.\ current velocity]
In a sliding-window estimator the trajectory window contains $n-1$ past
(already-recorded) poses and one current (active) pose $j=n$.  For past poses
$j<n$, $\bm{v}_{B,j}$ is a known, recorded quantity.  Only the current-pose
velocity $\bm{v}_{B,n}\equiv\bm{v}_{B}$ is the control input that the CBF can
design.  This distinction is critical for the barrier dynamics derived in
\Cref{sec:hdot}.
\end{remark}

% =========================================================================
\section{Bearing Vector Dynamics}
% =========================================================================

Differentiating \eqref{eq:feature_anchor_factored} with respect to time gives
\begin{align}
  {}^{C_j\to A}\dot{\bm{p}}_{f_i}
  &= {}^{A}\dot{\bm{p}}_{f_i}-{}^{A}\dot{\bm{p}}_{C_j} \\
  &= \dot{\rho}_{i,j}\bm{b}_{i,j}
     + \rho_{i,j}\dot{\bm{b}}_{i,j}. 
  \label{eq:relative_position_derivative}
\end{align}
The point feature is assumed to be static in the anchor frame, i.e.,
${}^{A}\dot{\bm{p}}_{f_i}=\bm{0}$.  Therefore,
\begin{equation}
  -{}^{A}\dot{\bm{p}}_{C_j}
  = \dot{\rho}_{i,j}\bm{b}_{i,j}
    + \rho_{i,j}\dot{\bm{b}}_{i,j}.
  \label{eq:relative_derivative_static_feature}
\end{equation}

\subsection{Depth Rate}

Using $\bm{b}_{i,j}=\eta_{i,j}\hat{\bm{b}}_{i,j}$ and multiplying
\eqref{eq:relative_derivative_static_feature} by
$\hat{\bm{b}}_{i,j}^{\top}/\eta_{i,j}$ gives
\begin{equation}
  \dot{\rho}_{i,j}
  = -\frac{\hat{\bm{b}}_{i,j}^{\top}}{\eta_{i,j}}
  {}^{A}\dot{\bm{p}}_{C_j}.
\end{equation}
Substituting \eqref{eq:pose_velocity} yields
\begin{resultbox}
\begin{equation}
  \boxed{
  \dot{\rho}_{i,j}
  = -\frac{\hat{\bm{b}}_{i,j}^{\top}}{\eta_{i,j}}
  \bm{R}_{C_j}^{A}\bm{R}_{B}^{C}\bm{v}_{B,j}}
  \label{eq:rhodot_multi_feature}
\end{equation}
\end{resultbox}

\subsection{Bearing Rate}

Substituting the depth rate back into
\eqref{eq:relative_derivative_static_feature} gives
\begin{align}
  \dot{\bm{b}}_{i,j}
  &= \frac{1}{\rho_{i,j}\eta_{i,j}^{2}}
     \bm{b}_{i,j}\bm{b}_{i,j}^{\top}
     \bm{R}_{C_j}^{A}\bm{R}_{B}^{C}\bm{v}_{B,j}
     -\frac{1}{\rho_{i,j}}\bm{R}_{C_j}^{A}\bm{R}_{B}^{C}\bm{v}_{B,j} \\
  &= -\frac{1}{\rho_{i,j}\eta_{i,j}^{2}}
     \left(\eta_{i,j}^{2}\bm{I}
     -\bm{b}_{i,j}\bm{b}_{i,j}^{\top}\right)
     \bm{R}_{C_j}^{A}\bm{R}_{B}^{C}\bm{v}_{B,j}.
\end{align}
Define
\begin{equation}
  \bm{\pi}_{i,j}
  = \eta_{i,j}^{2}\bm{I}
    -\bm{b}_{i,j}\bm{b}_{i,j}^{\top}.
  \label{eq:pij_def}
\end{equation}
Then the bearing dynamics become
\begin{keyresult}
\begin{equation}
  \boxed{
  \dot{\bm{b}}_{i,j}
  = -\frac{1}{\rho_{i,j}\eta_{i,j}^{2}}
  \bm{\pi}_{i,j}\bm{R}_{C_j}^{A}\bm{R}_{B}^{C}\bm{v}_{B,j}}
  \label{eq:bdot_multi_feature}
\end{equation}
\end{keyresult}

\begin{remark}
Because $\bm{b}_{i,j}$ is expressed in the anchor frame, the derived log-det
barrier depends directly on translational motion.  The pose dependence is
retained through $\bm{R}_{C_j}^{A}$ and ${}^{A}\bm{p}_{C_j}$.
\end{remark}

% =========================================================================
\section{Triangulation Information Matrix}
% =========================================================================

\subsection{Per-feature normal equations}

For feature $i$, the cross-product constraint at pose $j$ is
\begin{equation}
  \sx{\bm{b}_{i,j}}
  \left({}^{A}\bm{p}_{f_i}-{}^{A}\bm{p}_{C_j}\right)=\bm{0}.
\end{equation}
Stacking all $n$ constraints over the camera poses in the trajectory window
gives
\begin{equation}
  \underbrace{
  \begin{bmatrix}
  \bm{N}_{i,1}\\
  \vdots\\
  \bm{N}_{i,n}
  \end{bmatrix}}_{\bm{A}_i\in\R^{3n\times3}}
  {}^{A}\bm{p}_{f_i}
  =
  \underbrace{
  \begin{bmatrix}
  \bm{N}_{i,1}{}^{A}\bm{p}_{C_1}\\
  \vdots\\
  \bm{N}_{i,n}{}^{A}\bm{p}_{C_n}
  \end{bmatrix}}_{\bm{q}_i\in\R^{3n}},
  \qquad
  \bm{N}_{i,j}=\sx{\bm{b}_{i,j}}.
\end{equation}
The least-squares solution satisfies
$\bm{A}_{i}^{\top}\bm{A}_{i}{}^{A}\bm{p}_{f_i}
=\bm{A}_{i}^{\top}\bm{q}_{i}$.  Using
\begin{equation}
  \bm{N}_{i,j}^{\top}\bm{N}_{i,j}
  = \sx{\bm{b}_{i,j}}^{\top}\sx{\bm{b}_{i,j}}
  = \eta_{i,j}^{2}\bm{I}
    -\bm{b}_{i,j}\bm{b}_{i,j}^{\top}
  = \bm{\pi}_{i,j},
\end{equation}
the per-feature triangulation information matrix is
\begin{resultbox}
\begin{equation}
  \boxed{
  \bm{M}_{i}
  = \bm{A}_{i}^{\top}\bm{A}_{i}
  = \sum_{j=1}^{n}\bm{\pi}_{i,j}
  = \sum_{j=1}^{n}\left(\eta_{i,j}^{2}\bm{I}
  -\bm{b}_{i,j}\bm{b}_{i,j}^{\top}\right)
  \in\R^{3\times3}.}
  \label{eq:Mi_def}
\end{equation}
\end{resultbox}

\begin{remark}
$\bm{M}_{i}$ is positive semidefinite by construction.  It is positive
definite if the bearing set $\{\bm{b}_{i,j}\}_{j=1}^{n}$ provides sufficient
parallax for feature $i$.  Therefore, each feature has its own observability
metric and its own rank-deficiency risk.
\end{remark}

% =========================================================================
\section{Average Log-Determinant Barrier Function}
% =========================================================================

\subsection{Feature-wise and averaged observability metrics}

The feature-wise log-det metric is defined as
\begin{equation}
  \ell_i = \log\det\bm{M}_{i},
  \qquad i=1,\ldots,s.
  \label{eq:ell_i_def}
\end{equation}
The CBF does not read a single-feature metric.  Instead, it reads the average
log-det value over all selected features:
\begin{resultbox}
\begin{equation}
  \boxed{
  \bar{\ell}
  = \frac{1}{s}\sum_{i=1}^{s}\ell_i
  = \frac{1}{s}\sum_{i=1}^{s}\log\det\bm{M}_{i}.}
  \label{eq:average_logdet}
\end{equation}
\end{resultbox}
This average is the logarithm of the geometric mean of the determinants:
\begin{equation}
  \bar{\ell}
  = \log\left(\prod_{i=1}^{s}\det\bm{M}_{i}\right)^{1/s}.
\end{equation}

\subsection{Barrier function and safe set}

Let $h_{\min}\in\R$ denote the prescribed minimum acceptable average log-det
value.  Define
\begin{definition}[Average log-det barrier]
\begin{equation}
  h(\bm{x})
  = \bar{\ell}(\bm{x}) - h_{\min}
  = \frac{1}{s}\sum_{i=1}^{s}\log\det\bm{M}_{i}(\bm{x}) - h_{\min}.
  \label{eq:h_average_def}
\end{equation}
The safe set is
\begin{equation}
  \mathcal{C}
  = \left\{\bm{x}: h(\bm{x})\ge0\right\}
  = \left\{\bm{x}: \frac{1}{s}\sum_{i=1}^{s}\log\det\bm{M}_{i}(\bm{x})
  \ge h_{\min}\right\}.
  \label{eq:safe_set_average}
\end{equation}
\end{definition}

\begin{remark}
The average log-det condition enforces an aggregate observability requirement.
It does not, by itself, guarantee that every individual feature satisfies
$\log\det\bm{M}_{i}\ge h_{\min}$.  If feature-wise guarantees are required,
one should impose either $s$ individual CBF constraints or use
$\min_i\ell_i$-type logic.  In the present formulation, the requested CBF
constraint is based on the average metric $\bar{\ell}$.
\end{remark}

% =========================================================================
\section{Time Derivative of the Average Barrier Function}
\label{sec:hdot}
% =========================================================================

\subsection{Rate of change of each information matrix}

From \eqref{eq:Mi_def},
\begin{align}
  \dot{\bm{M}}_{i}
  &= \sum_{j=1}^{n}\dot{\bm{\pi}}_{i,j} \\
  &= -\sum_{j=1}^{n}
  \left(\dot{\bm{b}}_{i,j}\bm{b}_{i,j}^{\top}
  + \bm{b}_{i,j}\dot{\bm{b}}_{i,j}^{\top}\right),
  \label{eq:Midot}
\end{align}
where $\eta_{i,j}$ is assumed constant for normalized bearing measurements.
For unit-normalized bearings, $\eta_{i,j}=1$.

\subsection{Derivative of each feature-wise log-det metric}

For each positive-definite $\bm{M}_{i}$,
\begin{equation}
  \dot{\ell}_{i}
  = \frac{d}{dt}\log\det\bm{M}_{i}
  = \Tr\left(\bm{M}_{i}^{-1}\dot{\bm{M}}_{i}\right).
\end{equation}
Substituting \eqref{eq:Midot}, using cyclic trace identities, and exploiting
symmetry of $\bm{M}_{i}^{-1}$ gives
\begin{resultbox}
\begin{equation}
  \boxed{
  \dot{\ell}_{i}
  = -2\sum_{j=1}^{n}
  \bm{b}_{i,j}^{\top}\bm{M}_{i}^{-1}\dot{\bm{b}}_{i,j}.}
  \label{eq:ellidot_compact}
\end{equation}
\end{resultbox}

\subsection{Derivative of the average log-det barrier:
drift--control decomposition}

Since
$h=\frac{1}{s}\sum_{i=1}^{s}\ell_i-h_{\min}$, the barrier derivative is
\begin{equation}
  \dot{h}
  = \frac{1}{s}\sum_{i=1}^{s}\dot{\ell}_{i}.
\end{equation}
Substituting \eqref{eq:bdot_multi_feature} into
\eqref{eq:ellidot_compact} and noting that each pose carries its own velocity
$\bm{v}_{B,j}$ yields
\begin{align}
  \dot{h}
  &= \frac{2}{s}\sum_{i=1}^{s}\sum_{j=1}^{n}
  \frac{1}{\rho_{i,j}\eta_{i,j}^{2}}
  \bm{b}_{i,j}^{\top}\bm{M}_{i}^{-1}\bm{\pi}_{i,j}
  \bm{R}_{C_j}^{A}\bm{R}_{B}^{C}\bm{v}_{B,j}.
  \label{eq:hdot_average_expanded}
\end{align}
Because only the current-pose velocity $\bm{v}_{B,n}\equiv\bm{v}_{B}$ is the
control input while all past-pose velocities $\bm{v}_{B,j}$, $j<n$, are known
recorded quantities, the double sum splits naturally into a \emph{drift} term
and a \emph{control} term.  Using the symmetry identity
$\bm{b}_{i,j}^{\top}\bm{M}_{i}^{-1}\bm{\pi}_{i,j}
=\bigl(\bm{\pi}_{i,j}\bm{M}_{i}^{-1}\bm{b}_{i,j}\bigr)^{\top}$
to factor out the current-pose velocity:
\begin{keyresult}
\begin{equation}
  \boxed{
  \dot{h}
  = \underbrace{f(\bm{x})}_{\text{drift (past poses)}}
  + \underbrace{\bm{g}(\bm{x})^{\top}\bm{v}_{B}}_{\text{control (current pose)}}}
  \label{eq:hdot_drift_control}
\end{equation}
\end{keyresult}
where
\begin{align}
  f(\bm{x})
  &= \frac{2}{s}\sum_{i=1}^{s}\sum_{j=1}^{n-1}
  \frac{1}{\rho_{i,j}\eta_{i,j}^{2}}
  \bm{b}_{i,j}^{\top}\bm{M}_{i}^{-1}\bm{\pi}_{i,j}
  \bm{R}_{C_j}^{A}\bm{R}_{B}^{C}\bm{v}_{B,j},
  \label{eq:drift_def}\\[6pt]
  \bm{g}(\bm{x})
  &= \frac{2}{s}\sum_{i=1}^{s}
  \frac{1}{\rho_{i,n}\eta_{i,n}^{2}}
  \left(\bm{R}_{C_n}^{A}\bm{R}_{B}^{C}\right)^{\top}
  \bm{\pi}_{i,n}\bm{M}_{i}^{-1}\bm{b}_{i,n}
  \;\in\R^3.
  \label{eq:g_def}
\end{align}

\begin{remark}[Physical interpretation of the decomposition]
The drift $f(\bm{x})$ captures how \emph{past} translational motion has already
contributed to (or degraded) the observability metric.  It is fully determined
by the recorded trajectory and bearing history.  The control-input gradient
$\bm{g}(\bm{x})$ captures only the \emph{current} pose's influence, which is
the only quantity the controller can modify at the present time step.
\end{remark}

% =========================================================================
\section{Drift and Control-Input Gradient}
% =========================================================================

From the decomposition in \eqref{eq:hdot_drift_control}, the barrier
derivative has the control-affine structure
$\dot{h}=f(\bm{x})+\bm{g}(\bm{x})^{\top}\bm{v}_{B}$ with drift and
control-input gradient defined in \eqref{eq:drift_def}--\eqref{eq:g_def}.

\subsection{Drift term (past poses)}

The scalar drift collects the \emph{known} contributions of all past-pose
recorded velocities:
\begin{resultbox}
\begin{equation}
  \boxed{
  f(\bm{x})
  = \frac{2}{s}\sum_{i=1}^{s}\sum_{j=1}^{n-1}
  \frac{1}{\rho_{i,j}\eta_{i,j}^{2}}
  \bm{b}_{i,j}^{\top}\bm{M}_{i}^{-1}\bm{\pi}_{i,j}
  \bm{R}_{C_j}^{A}\bm{R}_{B}^{C}\bm{v}_{B,j}
  \;\in\R.}
  \label{eq:drift_boxed}
\end{equation}
\end{resultbox}
All quantities in $f(\bm{x})$ are evaluated from the recorded trajectory
window and are therefore computable at each time step without any design
freedom.

\subsection{Control-input gradient (current pose)}

The body-frame gradient that multiplies the designable velocity is
\begin{keyresult}
\begin{equation}
  \boxed{
  \bm{g}(\bm{x})
  = \frac{2}{s}\sum_{i=1}^{s}
  \frac{1}{\rho_{i,n}\eta_{i,n}^{2}}
  \left(\bm{R}_{C_n}^{A}\bm{R}_{B}^{C}\right)^{\top}
  \bm{\pi}_{i,n}\bm{M}_{i}^{-1}\bm{b}_{i,n}
  \;\in\R^3.}
  \label{eq:g_boxed}
\end{equation}
\end{keyresult}
The barrier derivative is therefore
\begin{keyresult}
\begin{equation}
  \boxed{
  \dot{h}(\bm{x})=f(\bm{x})+\bm{g}(\bm{x})^{\top}\bm{v}_{B}.}
  \label{eq:hdot_linear_average}
\end{equation}
\end{keyresult}

\begin{theorem}[Relative degree one]
The average log-det barrier in \eqref{eq:h_average_def} has relative degree
one with respect to the linear velocity input $\bm{v}_{B}$, since its first
time derivative depends affinely on $\bm{v}_{B}$.
\end{theorem}

\begin{remark}[Geometric interpretation]
In $\bm{g}(\bm{x})$, the term $\bm{M}_{i}^{-1}\bm{b}_{i,n}$ amplifies poorly
covered directions for feature $i$ at the current pose,
$\bm{\pi}_{i,n}$ removes the unobservable radial component along the bearing,
$1/\rho_{i,n}$ weights closer features more strongly, and
$(\bm{R}_{C_n}^{A}\bm{R}_{B}^{C})^{\top}$ maps the anchor-frame gradient into
the body-frame velocity coordinates.  The drift $f(\bm{x})$ summarises how
past motion has already shaped the information matrix.
\end{remark}

% =========================================================================
\section{Control Barrier Function Condition}
\label{sec:cbf}
% =========================================================================

\begin{definition}[Control Barrier Function]
A continuously differentiable function $h$ is a Control Barrier Function for
the safe set $\mathcal{C}=\{\bm{x}:h(\bm{x})\ge0\}$ if there exists a
class-$\mathcal{K}$ function $\alpha$ such that
\begin{equation}
  \sup_{\bm{v}_{B}\in\mathcal{V}}
  \left[\dot{h}(\bm{x},\bm{v}_{B})+\alpha(h(\bm{x}))\right]\ge0.
\end{equation}
\end{definition}

Using the linear class-$\mathcal{K}$ function $\alpha(h)=\gamma h$, with
$\gamma>0$, and substituting $\dot{h}=f(\bm{x})+\bm{g}(\bm{x})^{\top}\bm{v}_{B}$
from \eqref{eq:hdot_linear_average}, the average log-det CBF condition becomes
\begin{keyresult}
\begin{equation}
  \boxed{
  f(\bm{x})+\bm{g}(\bm{x})^{\top}\bm{v}_{B}
  +\gamma\left(
  \frac{1}{s}\sum_{i=1}^{s}\log\det\bm{M}_{i}(\bm{x})-h_{\min}
  \right)
  \ge0.}
  \label{eq:cbf_condition_average}
\end{equation}
\end{keyresult}
Equivalently, it enforces
\begin{equation}
  f(\bm{x}) + \bm{g}(\bm{x})^{\top}\bm{v}_{B} + \gamma h \ge 0,
  \label{eq:cbf_condition_expanded}
\end{equation}
which explicitly separates the drift from past poses and the control input from
the current pose.

% =========================================================================
\section{CBF-QP Formulation}
% =========================================================================

At each time instant, the nominal VIO/trajectory-tracking velocity command
$\bm{v}_{B,\mathrm{nom}}$ is minimally modified to satisfy the average
observability constraint:
\begin{keyresult}
\begin{equation}
  \boxed{
  \begin{aligned}
  \bm{v}_{B}^{*}
  ={}&\arg\min_{\bm{v}_{B}\in\R^{3}}
  &&\frac{1}{2}\|\bm{v}_{B}-\bm{v}_{B,\mathrm{nom}}\|_{\bm{W}}^{2}\\
  &\text{subject to}
  && f(\bm{x}) + \bm{g}(\bm{x})^{\top}\bm{v}_{B} + \gamma h \ge 0,
  \qquad [\text{average log-det CBF}]\\
  &&&\bm{v}_{B,\min}\le\bm{v}_{B}\le\bm{v}_{B,\max},
  \qquad [\text{velocity limits}]
  \end{aligned}}
  \label{eq:cbf_qp_average}
\end{equation}
\end{keyresult}
where $\bm{W}\succ0$ is a weighting matrix and
$\|\bm{a}\|_{\bm{W}}^{2}=\bm{a}^{\top}\bm{W}\bm{a}$.

\begin{theorem}[Forward invariance]
If the QP in \eqref{eq:cbf_qp_average} is feasible and solved at each time,
then the average safe set $\mathcal{C}$ in \eqref{eq:safe_set_average} is
forward invariant.  In particular,
\begin{equation}
  h(0)\ge0\quad\Longrightarrow\quad h(t)\ge0,
  \qquad \forall t\ge0.
\end{equation}
\begin{proof}[Proof sketch]
The CBF condition enforces $\dot{h}\ge-\gamma h$.  By the comparison lemma,
$h(t)\ge h(0)e^{-\gamma t}\ge0$ whenever $h(0)\ge0$.
\end{proof}
\end{theorem}

% =========================================================================
\section{Analytical Solution for a Single CBF Constraint}
% =========================================================================

When velocity box constraints are inactive and $\bm{W}=\bm{I}$, the QP admits
a closed-form solution.  If
$f(\bm{x})+\bm{g}^{\top}\bm{v}_{B,\mathrm{nom}}+\gamma h\ge0$, the nominal
command is already safe.  Otherwise, the active constraint projection gives
\begin{keyresult}
\begin{equation}
  \boxed{
  \bm{v}_{B}^{*}
  =
  \begin{cases}
  \bm{v}_{B,\mathrm{nom}},
  & f(\bm{x}) + \bm{g}^{\top}\bm{v}_{B,\mathrm{nom}} + \gamma h
  \ge 0,\\[6pt]
  \bm{v}_{B,\mathrm{nom}}
  -\dfrac{f(\bm{x}) + \bm{g}^{\top}\bm{v}_{B,\mathrm{nom}} + \gamma h}
  {\|\bm{g}\|_{2}^{2}}\bm{g},
  & f(\bm{x}) + \bm{g}^{\top}\bm{v}_{B,\mathrm{nom}} + \gamma h
  < 0.
  \end{cases}}
  \label{eq:closed_form_average}
\end{equation}
\end{keyresult}
The correction is the minimum-norm projection of the nominal command onto the
feasible half-space defined by the average log-det CBF constraint.

% =========================================================================
\section{Handling Unknown Depths}
% =========================================================================

The control-input gradient $\bm{g}$ in \eqref{eq:g_boxed} and the drift
$f(\bm{x})$ in \eqref{eq:drift_boxed} both depend on
$\rho_{i,j}$.  Two practical strategies are available:
\begin{table}[h]
\centering
\renewcommand{\arraystretch}{1.35}
\begin{tabular}{lll}
\toprule
\textbf{Strategy} & \textbf{Substitution} & \textbf{Trade-off} \\
\midrule
Online estimate
& $1/\rho_{i,j}\leftarrow1/\hat{\rho}_{i,j}$
& More accurate; requires depth initialization \\
Robust lower-bound model
& $1/\rho_{i,j}\leftarrow1/\rho_{\min}$
& Conservative; avoids explicit depth estimation \\
\bottomrule
\end{tabular}
\caption{Strategies for handling unknown feature depths.}
\label{tab:depth_multi}
\end{table}

% =========================================================================
\section{Complete Real-Time Algorithm}
% =========================================================================

\begin{algorithm}[H]
\caption{CBF-QP Safety Filter Using the Average Multi-Feature Log-Det Metric}
\label{alg:cbf_average}
\begin{algorithmic}[1]
\Require Bearings $\{\bm{b}_{i,j}\}_{i=1,j=1}^{s,n}$, depths or estimates
$\{\hat{\rho}_{i,j}\}_{i=1,j=1}^{s,n}$, rotations
$\{\bm{R}_{C_j}^{A}\}_{j=1}^{n}$, recorded past velocities
$\{\bm{v}_{B,j}\}_{j=1}^{n-1}$, nominal velocity
$\bm{v}_{B,\mathrm{nom}}$, parameters $h_{\min}$ and $\gamma$
\Ensure Safe linear velocity $\bm{v}_{B}^{*}$
\For{$i=1,\ldots,s$}
  \State $\bm{M}_{i}\leftarrow\sum_{j=1}^{n}\bm{\pi}_{i,j}$
  \Comment{Per-feature information matrix}
  \State $\ell_i\leftarrow\log\det\bm{M}_{i}$
  \Comment{Feature-wise log-det metric}
\EndFor
\State $\bar{\ell}\leftarrow\frac{1}{s}\sum_{i=1}^{s}\ell_i$
\Comment{Average observability metric}
\State $h\leftarrow\bar{\ell}-h_{\min}$
\State \textbf{--- Drift from past poses ---}
\State $f\leftarrow
\frac{2}{s}\sum_{i=1}^{s}\sum_{j=1}^{n-1}
\frac{1}{\hat{\rho}_{i,j}\eta_{i,j}^{2}}
\bm{b}_{i,j}^{\top}\bm{M}_{i}^{-1}\bm{\pi}_{i,j}
\bm{R}_{C_j}^{A}\bm{R}_{B}^{C}\bm{v}_{B,j}$
\State \textbf{--- Control gradient from current pose ---}
\State $\bm{g}\leftarrow
\frac{2}{s}\sum_{i=1}^{s}
\frac{1}{\hat{\rho}_{i,n}\eta_{i,n}^{2}}
(\bm{R}_{C_n}^{A}\bm{R}_{B}^{C})^{\top}
\bm{\pi}_{i,n}\bm{M}_{i}^{-1}\bm{b}_{i,n}$
\State $c_{\mathrm{safe}}\leftarrow-\gamma h - f$
\If{$\bm{g}^{\top}\bm{v}_{B,\mathrm{nom}}\ge c_{\mathrm{safe}}$}
  \State $\bm{v}_{B}^{*}\leftarrow\bm{v}_{B,\mathrm{nom}}$
\Else
  \State $\lambda^{*}\leftarrow
  \dfrac{c_{\mathrm{safe}}-\bm{g}^{\top}\bm{v}_{B,\mathrm{nom}}}
  {\|\bm{g}\|_{2}^{2}}$
  \State $\bm{v}_{B}^{*}\leftarrow\bm{v}_{B,\mathrm{nom}}+\lambda^{*}\bm{g}$
\EndIf
\State \Return $\bm{v}_{B}^{*}$
\end{algorithmic}
\end{algorithm}

% =========================================================================
\section{Summary of Key Equations}
% =========================================================================

\begin{table}[H]
\centering
\renewcommand{\arraystretch}{1.6}
\begin{tabular}{lll}
\toprule
\textbf{Quantity} & \textbf{Expression} & \textbf{Reference} \\
\midrule
Bearing of feature $i$ from pose $j$
& ${}^{C_j\to A}\bm{p}_{f_i}=\rho_{i,j}\bm{b}_{i,j}$
& \eqref{eq:feature_anchor_factored} \\
Pose velocity
& ${}^{A}\dot{\bm{p}}_{C_j}=\bm{R}_{C_j}^{A}\bm{R}_{B}^{C}\bm{v}_{B,j}$
& \eqref{eq:pose_velocity} \\
Bearing dynamics
& $\dot{\bm{b}}_{i,j}=-\frac{1}{\rho_{i,j}\eta_{i,j}^{2}}\bm{\pi}_{i,j}\bm{R}_{C_j}^{A}\bm{R}_{B}^{C}\bm{v}_{B,j}$
& \eqref{eq:bdot_multi_feature} \\
Projection matrix
& $\bm{\pi}_{i,j}=\eta_{i,j}^{2}\bm{I}-\bm{b}_{i,j}\bm{b}_{i,j}^{\top}$
& \eqref{eq:pij_def} \\
Per-feature information matrix
& $\bm{M}_{i}=\sum_{j=1}^{n}\bm{\pi}_{i,j}$
& \eqref{eq:Mi_def} \\
Barrier function
& $h=\bar{\ell}-h_{\min}$
& \eqref{eq:h_average_def} \\
Drift (past poses)
& $f(\bm{x})=\frac{2}{s}\sum_{i}\sum_{j=1}^{n-1}\frac{1}{\rho_{i,j}\eta_{i,j}^{2}}\bm{b}_{i,j}^{\top}\bm{M}_{i}^{-1}\bm{\pi}_{i,j}\bm{R}_{C_j}^{A}\bm{R}_{B}^{C}\bm{v}_{B,j}$
& \eqref{eq:drift_def} \\
Control gradient (current pose)
& $\bm{g}(\bm{x})=\frac{2}{s}\sum_{i}\frac{1}{\rho_{i,n}\eta_{i,n}^{2}}(\bm{R}_{C_n}^{A}\bm{R}_{B}^{C})^{\top}\bm{\pi}_{i,n}\bm{M}_{i}^{-1}\bm{b}_{i,n}$
& \eqref{eq:g_def} \\
Lie derivative
& $\dot{h}=f(\bm{x})+\bm{g}(\bm{x})^{\top}\bm{v}_{B}$
& \eqref{eq:hdot_linear_average} \\
CBF condition
& $f(\bm{x})+\bm{g}^{\top}\bm{v}_{B}+\gamma(\bar{\ell}-h_{\min})\ge0$
& \eqref{eq:cbf_condition_average} \\
\bottomrule
\end{tabular}
\caption{Summary of the drift--control multi-feature average log-det CBF derivation.}
\label{tab:summary_multi}
\end{table}

% =========================================================================
\section{Discussion}
% =========================================================================

\subsection{Meaning of the average log-det threshold}

The condition
\begin{equation}
  \bar{\ell}\ge h_{\min}
\end{equation}
means that the average feature-wise triangulation quality over the selected
feature set is kept above the desired level.  Equivalently, it enforces a lower
bound on the geometric mean of the determinants:
\begin{equation}
  \left(\prod_{i=1}^{s}\det\bm{M}_{i}\right)^{1/s}
  \ge \exp(h_{\min}).
\end{equation}
Thus, the CBF encourages commanded translations that improve the collective
observability of the features over the trajectory window.

\subsection{Comparison with individual-feature constraints}

The average formulation gives a single scalar CBF constraint and therefore
keeps the QP simple.  However, a strong feature may compensate for a weak one
inside the average.  If this is not acceptable, the QP can be modified with
per-feature drift--control constraints:
\begin{equation}
  f_i(\bm{x})+\bm{g}_{i}(\bm{x})^{\top}\bm{v}_{B}
  +\gamma\left(\ell_i-h_{i,\min}\right)\ge0,
  \qquad i=1,\ldots,s,
\end{equation}
where the per-feature drift and control gradient are
\begin{align}
  f_i(\bm{x})
  &= 2\sum_{j=1}^{n-1}
  \frac{1}{\rho_{i,j}\eta_{i,j}^{2}}
  \bm{b}_{i,j}^{\top}\bm{M}_{i}^{-1}\bm{\pi}_{i,j}
  \bm{R}_{C_j}^{A}\bm{R}_{B}^{C}\bm{v}_{B,j},\\[4pt]
  \bm{g}_{i}(\bm{x})
  &= \frac{2}{\rho_{i,n}\eta_{i,n}^{2}}
  \left(\bm{R}_{C_n}^{A}\bm{R}_{B}^{C}\right)^{\top}
  \bm{\pi}_{i,n}\bm{M}_{i}^{-1}\bm{b}_{i,n}.
\end{align}
The implemented formulation uses the average:
$f(\bm{x})=\frac{1}{s}\sum_{i}f_i$ and
$\bm{g}(\bm{x})=\frac{1}{s}\sum_{i}\bm{g}_i$.

\subsection{Connection to active perception}

The vector $\bm{g}(\bm{x})$ is the steepest first-order ascent direction of the
average log-det metric \emph{at the current pose} in the body-frame velocity
command space.  Therefore, the CBF-QP correction nudges the nominal velocity
toward translational motions that increase the average triangulation
observability while preserving nominal trajectory tracking as much as possible.
The drift $f(\bm{x})$ acts as a bias: if past motion has already provided
strong observability growth, the CBF constraint is easier to satisfy and less
correction is needed.

% =========================================================================
\appendix
\section{Skew-Symmetric Matrix Identity}
\label{app:skew_identity}

For any $\bm{a}\in\R^3$,
\begin{equation}
  \sx{\bm{a}}^2 = \bm{a}\bm{a}^{\top}-\|\bm{a}\|_2^2\bm{I}.
\end{equation}
Since $\sx{\bm{a}}^{\top}=-\sx{\bm{a}}$, it follows that
\begin{equation}
  \sx{\bm{a}}^{\top}\sx{\bm{a}}
  = -\sx{\bm{a}}^2
  = \|\bm{a}\|_2^2\bm{I}-\bm{a}\bm{a}^{\top}.
\end{equation}
With $\bm{a}=\bm{b}_{i,j}$, this gives the projection matrix
\begin{equation}
  \bm{\pi}_{i,j}
  = \eta_{i,j}^{2}\bm{I}-\bm{b}_{i,j}\bm{b}_{i,j}^{\top}.
\end{equation}

% =========================================================================
\begin{thebibliography}{9}

\bibitem{ames2019cbf}
A. D. Ames, S. Coogan, M. Egerstedt, G. Notomista, K. Sreenath, and P. Tabuada,
\textit{Control Barrier Functions: Theory and Applications},
Proc. European Control Conference (ECC), 2019.

\bibitem{pukelsheim2006}
F. Pukelsheim,
\textit{Optimal Design of Experiments},
SIAM, 2006.

\bibitem{hausman2017}
K. Hausman, J. Preiss, G. Sukhatme, and S. Weiss,
\textit{Observability-Aware Trajectory Optimization for Self-Calibration with Application to UAVs},
IEEE RA-L, 2017.

\bibitem{chaumette2006}
F. Chaumette and S. Hutchinson,
\textit{Visual Servo Control, Part I: Basic Approaches},
IEEE Robotics \& Automation Magazine, 2006.

\end{thebibliography}

\end{document}